% Source: https://artofproblemsolving.com/community/c3h3297239_number_theory_modulos
% Stable identity: aops-topic-3297239-q01
(a) Choose $u$ with $au\equiv1\pmod b$, possible by coprimality. Set $x=r+au(s-r)$. It is congruent to $r$ modulo $a$ and to $r+(s-r)=s$ modulo $b$. If $x,y$ are two solutions, both $a$ and $b$ divide $x-y$; because they are coprime, $ab$ divides it. Thus the solution class is unique.
(b) The system $x\equiv0\pmod6$, $x\equiv1\pmod{10}$, $x\equiv0\pmod{15}$ is impossible: its first condition makes $x$ even and the second makes it odd. Although the gcd of all three moduli is one, they are not pairwise coprime. This explains why the pairwise condition in the usual CRT matters.
